% ==== docspec-cas preamble (xelatex + CJK) ====
% docspec-cas.cls（改自 upstream cas-sc）已先載 [T1]fontenc + stix + inconsolata + booktabs/array/colortbl/
% dcolumn/xcolor/hyperref/makecell/multirow。此檔以 -H 在 class 之後注入，
% 只「覆蓋」字型與加 framed/fvextra/highlight，不重載 class 已有的套件。

% ---------- 版心：加寬(A4 寬頁 + 小邊距)，讓內文/表格更寬、左右留白少 ----------
% upstream cas-sc 原生 192mm 頁 / 13.7mm 邊距 (text 164.6mm) 偏窄留白多；覆寫成 A4 寬頁+小邊距。
\geometry{paperwidth=210mm,paperheight=297mm,hmargin=12mm,vmargin=20mm,headsep=12pt,footskip=14pt}

% ---------- 字型：fontspec + xeCJK 覆蓋 upstream cas-sc 的 stix/inconsolata ----------
\usepackage[numbers]{natbib}
\usepackage{fontspec}
\usepackage{xeCJK}
\XeTeXgenerateactualtext=1
\defaultfontfeatures{Ligatures=TeX}

\setmainfont{SourceSerif4}[Path=./fonts/, Extension=.otf,
  UprightFont=*-Regular, BoldFont=*-Bold, ItalicFont=*-It, BoldItalicFont=*-BoldIt]
\setsansfont{SourceSans3}[Path=./fonts/, Extension=.otf,
  UprightFont=*-Regular, BoldFont=*-Bold, ItalicFont=*-It, BoldItalicFont=*-BoldIt]
\setmonofont{SourceCodePro}[Path=./fonts/, Extension=.otf,
  UprightFont=*-Regular, BoldFont=*-Bold, ItalicFont=*-It, BoldItalicFont=*-BoldIt, Scale=0.90]

% CJK：TW-Kai 全字庫正楷體（標楷體 OFL 替身，內文）/ 思源黑（sans、標題）
% TW-Kai 單一字重、無 Bold → 粗體用真正的思源黑 Bold（不靠 AutoFakeBold 假粗體把字弄糊）。
% Extension=.ttf 只套用到主字（TW-Kai）；粗體 fallback 是 .otf，故 BoldFont 直接寫全名含副檔名。
\setCJKmainfont{TW-Kai-98_1}[Path=./fonts/, Extension=.ttf,
  BoldFont=SourceHanSansTC-Bold.otf]
\setCJKsansfont{SourceHanSansTC-Regular}[Path=./fonts/, Extension=.otf, BoldFont=SourceHanSansTC-Bold]
\setCJKmonofont{SourceHanSansTC-Regular}[Path=./fonts/, Extension=.otf, BoldFont=SourceHanSansTC-Bold]

% 符號 fallback：Source 拉丁字型缺 ↔ ↪ → 等箭頭、ℹ ✔ ✘ 等雜項符號，交給思源宋體渲染。
% 把這些碼點宣告成 CJK class，再設 CJK fallback family。
% 涵蓋：U+2100–214F（字母式符號，含 ℹ U+2139）、U+2190–21FF（箭頭）、
%       U+2460–24FF（圈號）、U+2500–257F（製表）、U+25A0–25FF（幾何）、
%       U+2600–27BF（雜項符號/dingbats，含 ✔ ✘ ⚠ 等文件常用標記）。
\xeCJKDeclareCharClass{CJK}{"2100 -> "214F, "2190 -> "21FF, "2460 -> "24FF,
  "2500 -> "257F, "25A0 -> "25FF, "2600 -> "27BF}
\setCJKfallbackfamilyfont{\CJKrmdefault}{SourceHanSerifTC-Regular}[Path=./fonts/, Extension=.otf]

% ---------- 字型紀律：杜絕 upstream cas-sc 載入的 stix/lmodern 漏進「文字」 ----------
% docspec-cas.cls（沿用 upstream cas-sc）在本檔之前 \RequirePackage{stix}（全載，含 text 編碼）並由 pandoc 牽連
% lmodern 重映數學符號字型。本檔的 fontspec \setmainfont 此刻剛把 TU 編碼的
% \rmdefault/\sfdefault/\ttdefault 設成 Source 三族。把這三個「現值」凍存，於文件
% 開始時釘回——擋掉任何在本檔之後（如 \begin{document} 前）改動預設族的套件，確保
% 全文 \normalfont/\textrm/\texttt 都落在 Source。數學若有符號仍由 stix 提供；
% 無數學式時 stix 完全不碰文字。
\let\docspecRMdefault\rmdefault
\let\docspecSFdefault\sfdefault
\let\docspecTTdefault\ttdefault
\AtBeginDocument{%
  \let\rmdefault\docspecRMdefault%
  \let\sfdefault\docspecSFdefault%
  \let\ttdefault\docspecTTdefault%
  \normalfont%
}

% ---------- 字級放大：upstream cas-sc 是期刊密排（10pt class），這不是投稿、要正常可讀 ----------
% upstream cas-sc 用 \ExecuteOptions{...10pt...} 硬綁 size10.clo（\normalsize≈10pt、\small≈9pt），
% 且各級字 def 散在 size10.clo / upstream cas-common。此處在 class 之後「整組重定義」NFSS 字級
% 巨集到 12pt 等級（與 size12.clo 同表），每級字一併放大 \baselineskip（非只放大字、行高不動），
% 並以 \AtBeginDocument 釘住 \normalsize 為現行字級，擋掉 class 後續把它重設回小字。
% （內含 \@setfontsize \p@ \@plus 等 @-internal 巨集，故須 \makeatletter 包夾。）
% 基準字級＝14.5pt（原 12pt 偏小、整份放大 ×1.21；每級字一併放大 baselineskip）。
\makeatletter
\renewcommand\normalsize{\@setfontsize\normalsize{14.5}{18.3}%
  \abovedisplayskip 14\p@ \@plus3\p@ \@minus7\p@
  \abovedisplayshortskip \z@ \@plus3\p@
  \belowdisplayshortskip 7.5\p@ \@plus3.5\p@ \@minus3\p@
  \belowdisplayskip \abovedisplayskip}
\renewcommand\small{\@setfontsize\small{13.2}{16.4}%
  \abovedisplayskip 13\p@ \@plus3\p@ \@minus6\p@
  \abovedisplayshortskip \z@ \@plus3\p@
  \belowdisplayshortskip 7.5\p@ \@plus3.5\p@ \@minus3\p@
  \belowdisplayskip \abovedisplayskip}
\renewcommand\footnotesize{\@setfontsize\footnotesize{12.1}{14.5}%
  \abovedisplayskip 12\p@ \@plus2\p@ \@minus5\p@
  \abovedisplayshortskip \z@ \@plus3\p@
  \belowdisplayshortskip 7\p@ \@plus3\p@ \@minus3\p@
  \belowdisplayskip \abovedisplayskip}
\renewcommand\scriptsize{\@setfontsize\scriptsize{9.7}{11.5}}
\renewcommand\tiny{\@setfontsize\tiny{7.25}{8.6}}
\renewcommand\large{\@setfontsize\large{17.4}{21.8}}
\renewcommand\Large{\@setfontsize\Large{20.8}{26.6}}
\renewcommand\LARGE{\@setfontsize\LARGE{25}{30}}
\renewcommand\huge{\@setfontsize\huge{30}{36}}
\renewcommand\Huge{\@setfontsize\Huge{30}{36}}
\makeatother
\AtBeginDocument{\normalsize}

% ---------- 章節標題字級：放大到大於 14.5pt 內文（upstream cas-common 寫死 12/11/10.5pt，現比內文小）----------
% cas 的 \sectionfont 等是 plain \def（載入期定義、之後不再改），在此 \def 覆蓋即生效。
% 用顯式 \fontsize（cas 對 NFSS 字級巨集的重定義無效，見內文字級修正）。\section≈17pt、
% \subsection≈15pt、\subsubsection≈14pt。
% ★用 \sffamily（＝CJKsans＝cjk_heading，預設標楷 TW-Kai）讓章節標題＝標楷；\bfseries 由
% CJKsans 的 AutoFakeBold 合成同面粗體（不換成黑體）。拉丁節號走 sans bold；內文/emphasis 仍宋體。
\def\sectionfont{\sffamily\fontsize{17pt}{20pt}\bfseries}
\def\ssectionfont{\sffamily\fontsize{15pt}{18pt}\bfseries\selectfont}
\def\sssectionfont{\sffamily\fontsize{14pt}{16pt}\fontseries{b}\fontshape{it}\selectfont}

% ---------- 中文段落形狀（行距放鬆到 CJK 舒適區）----------
\setlength{\parindent}{2em}
\setlength{\parskip}{0pt}
% 整體 leading：內文行距再放鬆（CJK/英文都更舒適）。1.6 × 14.5/18.3 基線 → 視覺約 2.0 行距。
\linespread{1.6}

% ---------- 顏色（xcolor 已由 docspec-cas 載入，僅 \definecolor）----------
% GitHub 風表格用色：表頭/隔列淡底 #F6F8FA、細格線 #D0D7DE（淺灰）。
\definecolor{docspecZebra}{HTML}{F6F8FA}
\definecolor{docspecRule}{HTML}{D0D7DE}
\definecolor{docspecCodeBg}{HTML}{F5F6F8}
\definecolor{docspecCodeRule}{HTML}{D9DEE6}
% inline code（GitHub 風淺灰圓角底）：底色比 block 稍深一階、細灰框。
\definecolor{docspecInlineBg}{HTML}{EFF1F3}
\definecolor{docspecInlineRule}{HTML}{D7DBE0}
\setlength{\heavyrulewidth}{1.1pt}
\setlength{\lightrulewidth}{0.5pt}
\arrayrulecolor{docspecRule}

% ---------- code 框：framed（docspec-cas 未載），灰底 + 軟換行 ----------
\usepackage{framed}
\usepackage{tcolorbox}
% 註：mermaid 佔位框用基本 tcolorbox（不載 skins/breakable）——breakable 牽連 pdfcol、
% skins 牽連 tikzfill.image，二者 TinyTeX medium scheme 未裝。佔位框只裝單張圖源（短），
% 不需跨頁；標準框已夠醒目。要長框跨頁＝把 pdfcol 加進 setup 再開 breakable。
\usepackage{fvextra}
\usepackage{tabularx}
\usepackage{xltabular}
\usepackage{seqsplit}
\setlength{\fboxsep}{0pt}
\newenvironment{docspeccode}%
  {\setlength{\OuterFrameSep}{5pt}%
   \def\FrameCommand{\fboxsep=9pt\colorbox{docspecCodeBg}}%
   \MakeFramed{\advance\hsize-2\fboxsep \FrameRestore}}%
  {\endMakeFramed}
\usepackage{etoolbox}
% pandoc 只在文件「真的有語法高亮 code block」時才定義 Shaded 環境（連帶 \KeywordTok 等
% token 巨集）。純散文/無 code 的快照不會有 Shaded → 直接 \renewenvironment 會
%「Environment Shaded undefined」整份炸。故用 ifdef 守：有就接管成灰框、沒有就不碰。
\ifcsname Shaded\endcsname
  \renewenvironment{Shaded}{\begin{docspeccode}}{\end{docspeccode}}
\fi
\BeforeBeginEnvironment{verbatim}{\begin{docspeccode}\small}
\AfterEndEnvironment{verbatim}{\end{docspeccode}}
\fvset{breaklines=true,breakanywhere=true,fontsize=\small,%
  breakindent=2em,breakautoindent=true,%
  breaksymbolleft={\footnotesize\color{docspecCodeRule}$\hookrightarrow$\,},%
  breaksymbolindentleft=0pt,breaksymbolsepleft=0.5em}

% ---------- 語法高亮 palette（需 pandoc --highlight-style=tango）----------
\definecolor{tokKeyword}{HTML}{B5106F}
\definecolor{tokString}{HTML}{0B3CC1}
\definecolor{tokComment}{HTML}{8A8A8A}
\definecolor{tokNormal}{HTML}{1A1A1A}
% token 巨集（\KeywordTok 等）同樣只在文件有語法高亮 code block 時由 pandoc 定義。
% 無 code 的快照不定義 → \renewcommand 會「undefined」整份炸。用 ifdef 守整塊：只在
% \KeywordTok 已存在時才覆寫 palette（pandoc 的 token 巨集是同一組、有其一即有全部）。
\AtBeginDocument{%
  \ifcsname KeywordTok\endcsname%
  \renewcommand{\KeywordTok}[1]{\textcolor{tokKeyword}{\textbf{#1}}}%
  \renewcommand{\ControlFlowTok}[1]{\textcolor{tokKeyword}{\textbf{#1}}}%
  \renewcommand{\ImportTok}[1]{\textcolor{tokKeyword}{\textbf{#1}}}%
  \renewcommand{\OtherTok}[1]{\textcolor{tokKeyword}{#1}}%
  \renewcommand{\FunctionTok}[1]{\textcolor{tokKeyword}{#1}}%
  \renewcommand{\BuiltInTok}[1]{\textcolor{tokKeyword}{#1}}%
  \renewcommand{\StringTok}[1]{\textcolor{tokString}{#1}}%
  \renewcommand{\SpecialStringTok}[1]{\textcolor{tokString}{#1}}%
  \renewcommand{\VerbatimStringTok}[1]{\textcolor{tokString}{#1}}%
  \renewcommand{\CharTok}[1]{\textcolor{tokString}{#1}}%
  \renewcommand{\DecValTok}[1]{\textcolor{tokString}{#1}}%
  \renewcommand{\FloatTok}[1]{\textcolor{tokString}{#1}}%
  \renewcommand{\BaseNTok}[1]{\textcolor{tokString}{#1}}%
  \renewcommand{\ConstantTok}[1]{\textcolor{tokString}{#1}}%
  \renewcommand{\CommentTok}[1]{\textcolor{tokComment}{\textit{#1}}}%
  \renewcommand{\CommentVarTok}[1]{\textcolor{tokComment}{\textit{#1}}}%
  \renewcommand{\DataTypeTok}[1]{\textcolor{tokKeyword}{#1}}%
  \renewcommand{\OperatorTok}[1]{\textcolor{tokNormal}{#1}}%
  \renewcommand{\AttributeTok}[1]{\textcolor{tokNormal}{#1}}%
  \renewcommand{\NormalTok}[1]{\textcolor{tokNormal}{#1}}%
  \fi%
}

% ---------- inline code：GitHub 風灰底，且「可斷行不凸出右界」 ----------
% 舊版用 tcolorbox \tcbox 做圓角灰底，但 \tcbox 是不可斷行的盒：長路徑（fleet/sc/**/status）
% 整盒無法換行 → 凸出右界被裁。soul 的 \hl 又無法吃 \allowbreak/\seqsplit 等命令（會壞）。
% 解法：把一段 code 切成「以 / 和 _ 為界的小段」，每小段各自包一個淺灰底 \colorbox
% （分隔符黏在前段尾，故灰底視覺連續、像一條灰帶），段與段之間插 \allowbreak —— 於是長
% code 能斷在任一分隔點、絕不凸出右界，而短 code 看起來就是一塊連續灰底。切段在 lua 的
% Code handler 完成（見 docspec-tables.lua dspx_code_segments），此處只定義「單段灰底盒」。
\let\oldtexttt\texttt
\setlength{\fboxsep}{0pt}
% 單段：等寬小一階 + 淺灰底，左右 0 間距讓相鄰段無縫銜接成連續灰帶。
\newcommand{\dspxcodeseg}[1]{\colorbox{docspecInlineBg}{\ttfamily\small\strut #1}}
% 後備：未切段時（極少數路徑）也能顯示成灰底等寬。
\newcommand{\dspxinlinecode}[1]{\dspxcodeseg{#1}}
\renewcommand{\texttt}[1]{\dspxcodeseg{#1}}
% 表格內 code：窄欄需可斷行。沿用無底、可斷的小號等寬（lua 已注入 \allowbreak），不套灰底
% 以免在 X 欄裡盒與換行打架。
\newcommand{\dspxcellcode}[1]{{\small\ttfamily #1}}

% ---------- #2 表格 / code block 上下對稱留白（不緊貼前後段）----------
% 表格與 code 環境前後加 \addvspace，避免「表格下緣緊貼下一段」「code 緊貼上一段」。
% \addvspace 會吸收合併相鄰間距，不會疊加成過大空白。
\newlength{\dspxblockskip}\setlength{\dspxblockskip}{10pt plus 3pt minus 2pt}
\BeforeBeginEnvironment{docspeccode}{\par\addvspace{\dspxblockskip}}
\AfterEndEnvironment{docspeccode}{\par\addvspace{\dspxblockskip}\noindent}

% ---------- list 收緊 ----------
\usepackage{enumitem}
\setlist{leftmargin=2em,topsep=4pt,itemsep=2pt,parsep=0pt}

% ---------- TikZ：原生向量繪圖（流程圖/狀態機；取代 mermaid 原始碼 dump）----------
% mermaid 區塊現由 docspec-tables.lua 渲成「可見佔位框」；release skill 教 agent 把每張
% mermaid 翻成等價 TikZ（raw latex block），於是圖原生向量化、CJK node 正常、零 Chromium。
% upstream cas-sc 已載 xcolor、xeCJK 已在上方設好 → TikZ node 內 CJK 直接走 \setCJKmainfont、正常。
% 載入時機在本檔末段（fontspec/xeCJK 已設定後），不與 docspec-cas 既有套件衝突
% （tikz 不重定義 docspec-cas 用到的 booktabs/array/colortbl 等）。
\usepackage{tikz}
\usetikzlibrary{positioning,arrows.meta,fit,backgrounds,calc,shapes.geometric}
% TikZ node 內中文：node 預設用 \normalfont（已是 Source+TW-Kai）；無須額外設定。
% 給 agent 的便利樣式（release skill 會引用）：圓角流程框、狀態圈、子圖背景框。
\tikzset{
  dspxflow/.style={rectangle, rounded corners=2pt, draw=docspecRule, line width=0.5pt,
    fill=docspecZebra, align=center, inner sep=5pt, font=\small},
  dspxstate/.style={ellipse, draw=docspecRule, line width=0.5pt, fill=docspecZebra,
    align=center, inner sep=3pt, font=\small},
  dspxedge/.style={-{Stealth[length=2.2mm]}, draw=docspecRule, line width=0.5pt},
  dspxgroup/.style={rectangle, rounded corners=3pt, draw=docspecRule, line width=0.4pt,
    fill=docspecCodeBg, inner sep=8pt},
}

% ---------- 頁尾「Page X of N」總頁數修正（off-by-one）----------
% upstream cas-common.sty 定義 \lastpage、並以 \AtEndDocument{\csxdef{lastpage}{\thepage}} 把末頁
% 頁碼寫進 .aux；但末頁 shipout 晚於該 hook，故記下的 N 永遠少 1（末頁顯示「of N-1」）。
% 正解＝lastpage 套件的 \pageref{LastPage}（它在末頁 \AtVeryEndDocument 後放標籤、精準）。
% 不改 docspec-cas.cls/docspec-cas-common.sty 原檔：只在此覆蓋 cas 用來算 N 的巨集 \lastpage。
% 時機＝\AtBeginDocument 末段（after cas 讀 .aux 重設 \lastpage 之後），把 \lastpage 重綁
% 成 \pageref{LastPage}；頁尾 \__first_foot:/\__cas_foot: 在 shipout 取 \lastpage 即得正解。
\usepackage{lastpage}
\AtBeginDocument{\renewcommand{\lastpage}{\pageref{LastPage}}}

% ---------- strikethrough：pandoc 把 ~~text~~ 轉成 \st{...}（需 soul 套件）----------
% soul 對 CJK 相容性差（逐 byte 處理多位元字元）；改用 ulem 的 \sout（逐行繪刪除線）。
% [normalem] 保留 \emph 原本行為（不被 ulem 改成底線），只借用 \sout。
% \providecommand 防止 docspec-cas/docspec-cas-common 若已定義 \st 時衝突。
\usepackage[normalem]{ulem}
\providecommand{\st}[1]{\sout{#1}}

% ---------- 寡行控制 ----------
\clubpenalty=10000
\widowpenalty=10000
\displaywidowpenalty=10000
\brokenpenalty=10000
